
\documentclass[11pt, a4paper]{article}

% ===== ESSENTIAL PACKAGES =====
\usepackage{geometry}
\geometry{a4paper, margin=1.2in}
\usepackage{setspace}
\usepackage{array}
\usepackage{xcolor}
\onehalfspacing
\setlength{\parskip}{1em}
\usepackage{multicol}

\usepackage{catchfile}

% ===== WORD COUNT COMMAND =====
\newcommand{\wordcount}{%
    \ifnum\pdf@shellescape=1%
        \immediate\write18{texcount -1 -sum Oblig1.tex > wordcount.txt}%
        \CatchFileDef{\words}{wordcount.txt}{}%
        \words%
    \else%
        [Enable -shell-escape]%
    \fi%
}
\usepackage{mdframed}

\newmdenv[
    leftmargin=-50pt,
    rightmargin=0pt,
    innerleftmargin=3pt,
    innerrightmargin=0pt,
    innertopmargin=5pt,
    innerbottommargin=5pt,
    skipabove=5pt,
    skipbelow=5pt,
    linewidth=1pt,
    linecolor=black,
    topline=true,
    rightline=false,
    bottomline=false
]{sectionboxinner}

\newenvironment{sectionbox}[1][]{%
    \begin{sectionboxinner}
    \textbf{\large #1}
    \vspace{1pt}
    \newline
}{%
    \end{sectionboxinner}
}

\newmdenv[
    leftmargin=0pt,
    rightmargin=0pt,
    innerleftmargin=0pt,
    innerrightmargin=0pt,
    innertopmargin=3pt,
    innerbottommargin=3pt,
    skipabove=8pt,
    skipbelow=8pt,
    linewidth=1.2pt,
    linecolor=gray!60,  % Lighter gray color
    topline=false,
    rightline=false,
    bottomline=false
]{subsectionboxinner}

\newenvironment{subsectionbox}[1][]{%
    \begin{subsectionboxinner}
    \textbf{#1}
    \vspace{3pt}
    \newline
}{%
    \end{subsectionboxinner}
}

% ===== CUSTOM ENVIRONMENTS FOR PHILOSOPHICAL ARGUMENTS =====
\newenvironment{argument}{\begin{quote}\itshape\textbf{Argument: }}{\end{quote}}
\newenvironment{objection}{\begin{quote}\itshape\textbf{Innvending: }}{\end{quote}}
\newenvironment{reply}{\begin{quote}\itshape\textbf{Reply: }}{\end{quote}}

% ===== DOCUMENT BEGIN =====
\title{Kap. 14 - Hva krever moralen av oss? At vi gjør vår plikt}
\author{Oliver Ekeberg}
\date{\today}

\begin{document}
\maketitle

\tableofcontents


\begin{flushleft}
    Moralens grunnlag bør være p fremme lykke. dvs. det vi bør gjøre, er dermed det som fører til mest mulig lykke.

    \begin{quote}
        \textbf{\textit{Lykke:}} velbehag, nytelse, fravær av smerte
    \end{quote}

    En handling er moralsk om den øker totale lykken. Dette kan trekke flere pareleller til mikroøkonomi. Der beregner vi "nytte" i et marked og forsøker å maksimere det for å gjøre markedet perfekt.


    Jeg kan bruke 45 kroner på en god kopp med kaffe hver dag. Det vil gi meg lykke. Men hvis jeg går ned til svart kafffe som koster 20kr, kan jeg donere 25 til veldedighet for vaksinering av malaria. Det vil øke den totale lykken, antageligvis. 
\end{flushleft}


\begin{sectionbox}[Lykkemaksimering]
\begin{flushleft}
    \begin{quote}
        \textbf{\textit{Utilitarisme: størst mulig lykke prinsippet.}}
    \end{quote}
    
    Handlinger er riktige i den grad de fremmer lykke, gale i den grad de produserer det motsatte av lykke - smerte.\newline
    

    Dette prinsippet kalles også "nytteprinsippet". Prinsippet er normativt i den forstand at vi bør handle på måte som gjør at vi produserer så mye lykke vi kan. Under følger noen kjennetegn ved prinsippet

    \begin{itemize}
        \item Forutsetter at nytelser er homogent. dvs. lykke og nytelse er samme ting, slik at i prinsippet er mulig å sammenlikne, måle og telle nytelser
        \item Krever at man beregner med alle som berøres av en handling. Ikke bare en selv
        \item Det krever upartiskhet mellom egen og andres lykke
        \item Det vektlegger de mulige tilstandene som følger av handlingene
        \item Det er et maksimeringsprinsipp. Lykke-funksjonen er konkav.
    \end{itemize}

    Vi har nå presentert kjernen i \textit{utilitarisme}. Den kan forklare hvorfor enhver handling er rett eller gal. Den er også lett å forstå og dermed anvendelig. MEN den er faktisk vanskelig å anvende. Hvordan skal man måle lykke kvantitativt eller kvalitativt? En forutsetning er at lykke er målbart, eller kan telles. Dette viser seg å være en utfordring.
\end{flushleft}
\end{sectionbox}


\begin{sectionbox}[Regler og ulike typer lykke]
\begin{flushleft}
    Kan vi egentlig vite hva konsekvensene av en handling skal være? Den handlingen du vurderer å begå er jo en kausal effekt av alt som har skjedd opptil det punktet der. Og alt som skal skje kan påvirkes av det du gjør i det øyeblikket. Er kanskje min subjektive vurdering av sannsynlig helt ulik objektiv sannsynlighet? \newline

    Det er to distinksjoner du bør legge merke til: Mill (John Stuart Mill) mener man kan skille kvalitativt mellom ulike typer nytelse eller ubehag. Han argumenterer for at det finnes noen lykketilstander som er bedre enn andre. Lykke betyr ikke bare behagelighet. Det å trene kan virke ubehagelig for noen, men det er klart at uten trening, vil man dø tidligere, og leve et verre liv.  \newline

    En annen viktig distonksjon Mill antyder, er mellom det som senere er blitt kalt regel og handlingsutilitarisme.\newline

    Det er viktig å telle med både \textit{kvantitativ} og \textit{kvalitativ} nytte. For mye av noe er ikke alltid bra

    \begin{quote}
        \textbf{\textit{Regelutilitarisme: }} en totrinns modell: det som er grunnen til at noe er en moralsk regel, er at den er lykkemaksimerende. Dermed trenger du ikke regne på hvilken handling som tilfredsstiller størst mulig lykke
    \end{quote}
\end{flushleft}
\end{sectionbox}

\begin{sectionbox}[Kritisk potensial]
\begin{flushleft}
    Utilitarismen er forskjellig fra aristotelisk og kantiansk etikk, fordi i utilitarisme er vi ikke interessert i karakter, motivasjon eller følelser som ligger bak handlingen. Vi er istedenfor interessert i hvordan det går i verden når aktører handler. Teorien er derfor ikke ute etter å evaluere folk, men teorien er praktisk, rasjonell og fremtidsrettet.\newline

    En styrke med utilitarisme er at den vurderer konsekvenser av å handle i tråd med nedarvet normer eller et moralsk påbud faktisk er gode. Mill var en progressiv viktig tenker på 1800-tallet i Storbrittania. Han brukte utilitarisme i å argumentere for kvinners rettigheter, i tillegg til å utvide tilgang til utdanning og fritid til mennesker fra alle samfunnslag.
\end{flushleft}
\end{sectionbox}


\end{document}